\documentclass{article}

\usepackage{fontspec,xunicode,xltxtra,lipsum,amsmath}
\usepackage{titlesec}
\usepackage[top=1in,bottom=1in,left=1.25in,right=1.25in]{geometry}
\usepackage{xcolor}
\usepackage{xunicode}
\usepackage{xltxtra}
\usepackage{polyglossia}
\usepackage{verbatim}
\usepackage{verbatimbox}
\usepackage{graphicx} 		% 插入图片
\usepackage{setspace}
\usepackage{indentfirst}             									% 段首缩进
\XeTeXlinebreaklocale "zh"    	           							% 中文断行
\XeTeXlinebreakskip = 0pt plus 1pt minus 0.1pt			% 中文断行
\renewcommand{\baselinestretch}{1.25}

\defaultfontfeatures{Mapping=tex-text}

% \setdefaultlanguage{english}  		% 设置默认语言 english | chinese

% 设置字体
\setsansfont{Calibri}
\setmainfont[BoldFont=SimHei]{SimHei}

\titleformat{\section}{\Large\song}{\thesection}{1em}{}

% 新定义字体
\newfontfamily\shu{FZShuTi}                           % 方正舒体
%\newfontfamily\nming{MingLiU_HKSCS} 		% 新细明体
\newfontfamily\lishu{LiSu} 								% 隶书
\newfontfamily\hei{SimHei} 								% 黑体
\newfontfamily\yahei{Microsoft YaHei} 			% 微软雅黑
\newfontfamily\henghei{Microsoft JhengHei}  % 微軟正黑體
\newfontfamily\song{SimSun}  			% 宋体
\newfontfamily\nsong{NSimSun}  		% 新宋体
\newfontfamily\yao{FZYaoTi}                % 方正姚体
\newfontfamily\ming{MingLiU} 			% 细明体
\newfontfamily\youyuan{YouYuan}		% 幼圆
\newfontfamily\kai{KaiTi}         			% 楷体
\newfontfamily\cai{STCaiyun}   			% 华文彩云
\usepackage{enumitem}               % list
\setlist[enumerate]{itemsep=0mm}    % List Line space

\usepackage{paralist}

\usepackage{listings} % needed for the inclusion of source code
\lstset{basicstyle = \ttfamily, frame      = single}

% \makeindex                             % 加索引
% \index{法华经}

\begin{document}

\title{\hei XeTeX使用小结}
\author{\song 江长里}
\date{\kai 2018年2月21日}
\maketitle

\tableofcontents
\newpage


%%%% 段落首行缩进两个字 %%%%
\makeatletter
%\let\@afterindentfalse\@afterindenttrue
%\@afterindenttrue
%\makeatother
%\setlength{\parindent}{2em}  %中文缩进两个汉字位

\section{简介}
(默认字体) 以前使用CJK进行中文的排版，需要自己生成字体库，近日，出现了XeTeX，\yao （姚体）可以比较好的解决中文字体问题，不需要额外
生成LaTeX字体库，\kai （楷体）直接使用计算机 \henghei （微軟正黑體）系统里的字体。



\section{列出source code}
  
\lstinputlisting[language=Tex]{source.tex}

\section{List}
%\begin{itemize}
	\item 工人
	\item 农民 
	\item 学生 
\end{itemize}
\begin{compactitem}
	\item 《霍元甲》:  
	\item   《妙法莲华经》:  
	\item    寒山拾得: 
\end{enumerate}

\lstinputlisting[language=Tex]{list.tex}




\section{字体列表}
本文使用了大量本机自带的字体。
%%%%%%%%%%%%%%%%%%%%%%%%


%%%%%%%%%%%%%%%%%%%%%%%%%%%%
\begin{table}[htbp]
\caption{字体列表 office station}

\centering
\begin{tabular}{|l|c|r|}
\hline
\hei 字体 & \hei 命令 & \hei 字体效果 \\
\hline

\kai 宋体 & \verb+\song+ & \song 宋体 \\
\kai 新宋体 & \verb+\nsong+ & \nsong 新宋体 \\
\kai 楷体 & \verb+\kai+ & \kai 楷体 \\
\kai 黑体 & \verb+\hei+ & \hei 黑体 \\
\kai 隶书 & \verb+\lishu+ & \lishu 隶书 \\
\kai 方正舒体 & \verb+\shu+ & \shu  方正舒体 \\
\kai 方正姚体 & \verb+\yao+ & \yao 方正姚体 \\
\kai 幼圆 & \verb+\youyuan+ & \youyuan 幼圆 \\
\kai 微软雅黑 & \verb+\yahei+ & \yahei 微软雅黑 \\
\kai 微軟正黑體 & \verb+\henghei+ & \henghei 微軟正黑體 \\
\kai 細明體 & \verb+\ming+ & \ming 細明體 \\
\kai 新细明体 & \verb+\nming+ & \nming 新细明体 \\
\kai 华文彩云 & \verb+\cai+ & \cai 华文彩云 \\


\hline
\end{tabular}
\end{table}





图片插入效果图：
    \begin{center}
		\includegraphics[width=1\textwidth]{sea.jpg}
	\end{center}


% \printindex                             % 印索引	
	
\end{document}

